\section{Which orders are settled}\label{sec:further}

Each formula needs two facts at a given \(p\). The first is the
rigidity condition \((\mathrm R_{p-1})\) of Section~\ref{sec:counting}.
It gives \(N_p\), \(M_p\), \(F_p\) and \(2F_p\) as counts with
multiplicity: the multiplicities of the pairs add up to these numbers.
The second is that every rational-exponential pair is simple, that is,
\(\Xp\) is reduced. Then the pairs are distinct, and the elliptic pairs
are distinct on a generic lattice. On a special lattice several
elliptic pairs can merge, while the count with multiplicity stays the
same; Section~\ref{sec:ks} shows this at order three.
Table~\ref{tab:status} records what is known about each fact.

\begin{table}[htbp]
\centering
\caption{What is known. For elliptic pairs, distinctness is asserted
only on a generic lattice.}
\label{tab:status}
\begin{tabular}{@{}p{0.36\linewidth}p{0.58\linewidth}@{}}
\toprule
Statement&Known for\\
\midrule
Counts with multiplicity&\(2p-1\) a prime or a prime power (proved);
every \(p\le73\) (exact certificate); every \(p\le573\) (finite-field
computations)\\[2pt]
All \(N_p\) pairs distinct&\(2p-1\) prime (proved); \(p=5,8,11\) (exact
computations); hence every \(p\le12\)\\[2pt]
All \(F_p\) pairs of the pure subfamily distinct&whenever all \(N_p\)
pairs are distinct; every \(p\le14\) (exact computations)\\[2pt]
Both statements, every \(p\)&conjectured\\
\bottomrule
\end{tabular}
\end{table}

The prime-index argument excluded solutions at infinity by reducing
the continuous convolution modulo \(\ell\). The same argument proves
\eqref{eq:length} whenever \(2p-1=\ell^e\) is an odd prime power
(Appendix~\ref{app:verification}). It gives the count with
multiplicity, but it does not show that the pairs are distinct.

More generally, the exact arithmetic certificates supplied with the
paper establish \((\mathrm R_d)\) for every \(d\le72\), hence
\begin{equation}\label{eq:coverage}
 \len\Xp=\binom{2p-1}{p-1},\qquad
 \len\mathcal E_p(g_2,g_3)=2\binom{2p-2}{p-2}
 \qquad(2\le p\le73).
\end{equation}
In the same range, Theorem~\ref{thm:symmetric} and
Corollary~\ref{cor:elliptic-symmetric} give
\(\len\Xp^{\mathrm{sym}}=F_p\) and, for even \(p\),
\(\len\mathcal E_p^{\mathrm{sym}}(g_2,g_3)=2F_p\).
Together with the ancillary proof, the certificates also cover infinite
families of higher orders.
Finite-field Gr\"obner computations, checked by independent pipelines,
establish \((\mathrm R_d)\) for every \(d\le572\); they extend
\eqref{eq:coverage} and the two symmetric lengths, as counts with
multiplicity, to every \(2\le p\le573\) (repository, directory
\texttt{certificates/rigidity\_572/}).
The certificates hold for all complex coefficients, not only real or
rational ones. The ancillary proof and the full certificate are described
in Appendix~\ref{app:verification}.

Equation \eqref{eq:coverage} is a statement about counts with
multiplicity only. At composite indices, exact computations show that
all pairs are simple at \(p=5\), \(p=8\) and \(p=11\)
(Appendix~\ref{app:verification}). With Theorem~\ref{thm:prime}, all
pairs are therefore distinct for every \(p\le12\), and
Theorem~\ref{thm:elliptic-count} gives \(M_p\) distinct elliptic pairs
on a generic lattice at these orders. Inside the two pure subfamilies
the computation is much smaller, and it shows that the \(F_p\) pairs
are distinct for every \(p\le14\), which includes the composite indices
\(9\), \(15\), \(21\), \(25\) and \(27\).
